\color{unidarkgreen}{\subsection{Логика отключения смежного трансформатора}}
\color{black}

\begin{enumerate}

\item
Ввод в работу логики отключения смежного трансформатора осуществляется с помощью уставки \textbf{«ЛО смеж Т»}. Оперативный ввод/вывод ЛО см. Т от ТЗНП производится с помощью БУ (блока управления функциями) «Вывод ЛО смеж. Т». 

\item
С целью исключения возникновения недопустимого режима работы параллельного трансформатора с изолированной нейтралью при замыкании на землю одной фазы. Предусматривается возможность первым действием сформировать сигнал \textit{«Откл. смеж. Т»} на отключение выключателя стороны ВН смежного трансформатора. Логика отключения смежного трансформатора формируется по пуску ТЗНП ВН с независимой выдержкой времени, которая задается уставкой \textbf{«Тср ТЗНП откл смеж Т»}. 

\item
Алгоритм работы логики отключения смежного трансформатора приведен на рисунке \ref{pic:LO_T_TZNP}. Параметры для настройки логики отключения смежного трансформатора приведены в таблице \ref{app_set:losmtr}.

\medskip
\begin{figure}[H]
\centering
\includegraphics[width=1\textwidth,height=1\textheight,keepaspectratio]{\fbpath/06.01 ДЗТ/021.0361 ЛО смеж Т от ТЗНП/img/LO_T_TZNP.pdf}
\captionof{figure}{Функциональная схема логики отключения смежного Т}\label{pic:LO_T_TZNP}
\end{figure}


\end{enumerate}